% =====================================================================
%  Chapter 05 — Layers 1 and 2
% =====================================================================
\chapter{Layer 1 and Layer 2: The Property Engine and the Two-Dimensional Solvers}
\label{ch:layer12}

The numerical programme was built before the analytical programme, and it was
built bottom-up. This chapter covers the two lowest layers: the property
engine that supplies every thermophysical number in the system, and the
two-dimensional spectral solvers in which each mechanism was first isolated.
All figures quoted are from \texttt{results/results.db}.

\section{Layer 1: the property engine}
\label{sec:layer1}

\texttt{layer1/tfirst\_props.py} is a single module and a hard rule: no solver
file may inline a property formula. It exposes six entry points.

\begin{center}
\small
\begin{tabular}{@{}l p{0.62\textwidth}@{}}
\toprule
Function & What it returns \\
\midrule
\texttt{scalar\_T\_props(T)}
  & \(\mu(T)\), \(k(T)\), \(\rho\), \(c_v\) for the working fluid. \\
\texttt{A\_field(T)}
  & The diffusivity field \(A=k(T)/(\rho c_v)\). Asserts strict positivity
    pointwise on every call in debug mode. \\
\texttt{second\_law\_check(fluid)}
  & Must pass for a fluid before any solver run is permitted. \\
\texttt{route1\_coeffs(T)}
  & \(\mu_{\min},\mu_{\max},A_{\min},A_{\max}\): the uniform bounds a
    regularity margin would need. \\
\texttt{route2\_theta(u,\(\nu\))}
  & The auxiliary scalar source for Route~2. \\
\texttt{self\_similar\_props(\(\lambda\))}
  & Coefficients along a self-similar profile. \\
\bottomrule
\end{tabular}
\end{center}

The layer is small and carries \(71\) tests. Its importance is
disproportionate to its size for one reason: because every \(A\) value in the
programme comes from one code path, \(A\) values are comparable \emph{across
solvers, dimensions and years}. Chapter~\ref{ch:layer3} uses that
comparability to settle a real question --- whether \(A_{\min}\) in a violent
adversarial 3D run is set by the flow or by the background temperature --- by
comparing a number produced in Layer~3 against a number produced in Layer~2.

\section{Layer 2: the two-dimensional solvers}
\label{sec:layer2}

Seven modules, four milestones. \texttt{self\_similar\_IC.py} is the single
source of adversarial initial data. \texttt{T\_solver\_2D.py} is the Route-1
solver: incompressible Navier--Stokes with \(\mu(T)\) plus the scalar
transport equation for \(T\). \texttt{route2\_2D.py} is the Route-2 solver:
the exact Prize equations plus \(\theta\). \texttt{LPS\_monitor.py} and
\texttt{stretching\_2D.py} are diagnostics; \texttt{lambda\_sweep\_2D.py} and
\texttt{mu\_limit\_2D.py} are the two experiment drivers.

Two bugs found in \Sn{02} are worth recording because they are the kind that
produce plausible wrong numbers rather than crashes.

\begin{record}
\textbf{Dealiasing broke Hermitian symmetry.} The \(2/3\) filter zeroed the
row at \(k=+N/3\) but not the row at \(k=-N/3\), so the real-field FFT lost
its conjugate symmetry. Fixed by zeroing the closed range \(\abs{k}>N/3\).

\smallskip
\textbf{The energy formula went negative.} \(E=\tfrac12\langle\psi\omega\rangle
L^2\) is negative for dipole vortices, where \(\psi\) and \(\omega\) are
anti-correlated. Replaced by \(E=\tfrac12\langle\abs{\nabla\psi}^2\rangle L^2\),
computed spectrally, which is non-negative by construction.
\end{record}

\section{The \texorpdfstring{\(\lambda\)}{lambda} sweep (M3)}
\label{sec:lambda-sweep}

The sweep tests the conjecture that temperature-dependent viscosity defeats
self-similar blow-up for every \(\lambda>0\) --- that no critical threshold
\(\lambda_c\) exists. Ten runs, all at \(N=64\), \(t_{\mathrm{end}}=0.5\),
ideal-gas properties, logged under \texttt{claim\_id} \texttt{Conjecture\_3.5}
and its two variants.

\begin{table}[ht]
\centering
\small
\begin{tabular}{@{}l l r r r l@{}}
\toprule
\texttt{exp\_id} & profile & \(\lambda\) & \(A_{\min}\) & suppression &
verdict \\
\midrule
\texttt{EXP-L2-R1-CCF-001} & CCF & \(1.2000\) & \(3.0801\!\times\!10^{-5}\) & \(9.1896\) & \textsf{PASS} \\
\texttt{EXP-L2-R1-CCF-002} & CCF & \(0.9000\) & \(3.0801\!\times\!10^{-5}\) & \(7.4643\) & \textsf{PASS} \\
\texttt{EXP-L2-R1-CCF-003} & CCF & \(0.6057\) & \(3.0801\!\times\!10^{-5}\) & \(6.0869\) & \textsf{PASS} \\
\texttt{EXP-L2-R1-CCF-004} & CCF & \(0.4703\) & \(3.0801\!\times\!10^{-5}\) & \(5.5416\) & \textsf{PASS} \\
\texttt{EXP-L2-R1-CCF-005} & CCF & \(0.3000\) & \(3.0801\!\times\!10^{-5}\) & \(4.9246\) & \textsf{PASS} \\
\texttt{EXP-L2-R1-CCF-006} & CCF & \(0.1000\) & \(3.0801\!\times\!10^{-5}\) & \(4.2871\) & \textsf{PASS} \\
\texttt{EXP-L2-R1-BOUS-001} & Bouss. & \(1.9206\) & \(3.0801\!\times\!10^{-5}\) & \(15.143\) & \textsf{PASS} \\
\texttt{EXP-L2-R1-BOUS-002} & Bouss. & \(1.3991\) & \(3.0801\!\times\!10^{-5}\) & \(10.549\) & \textsf{PASS} \\
\texttt{EXP-L2-R1-BOUS-003} & Bouss. & \(1.1843\) & \(3.0801\!\times\!10^{-5}\) & \(9.0901\) & \textsf{PASS} \\
\texttt{EXP-L2-R1-ADV-001} & adv.\ min & \(0.0100\) & \(3.1740\!\times\!10^{-5}\) & \(4.0278\) & \textsf{PASS} \\
\bottomrule
\end{tabular}
\caption[The \(\lambda\) sweep]{The M3 \(\lambda\) sweep, all ten runs.
``Suppression'' is \texttt{suppression\_ratio\_final}, the ratio of the
viscous damping rate to the self-similar blow-up rate. The verdict recorded
for the conjecture was: no \(\lambda_c\) found; all \(\lambda\)
suppressed.}
\label{tab:lambda-sweep}
\end{table}

\begin{figure}[ht]
\centering
\begin{tikzpicture}
\begin{axis}[
  width=0.82\textwidth, height=6.2cm,
  xlabel={self-similar exponent \(\lambda\)},
  ylabel={suppression ratio at \(t=0.5\)},
  xmin=0, xmax=2.05, ymin=0, ymax=17,
  grid=both, grid style={gray!18},
  legend pos=north west, legend cell align=left,
  legend style={font=\scriptsize, draw=gray!40},
  tick label style={font=\scriptsize}, label style={font=\small}
]
\addplot[mark=*, thick, color=openblue] coordinates {
  (0.1,4.2871) (0.3,4.9246) (0.4703,5.5416) (0.6057,6.0869)
  (0.9,7.4643) (1.2,9.1896)
};
\addlegendentry{CCF profiles}
\addplot[mark=square*, thick, color=provedgreen] coordinates {
  (1.1843,9.0901) (1.3991,10.549) (1.9206,15.143)
};
\addlegendentry{Boussinesq profiles}
\addplot[mark=triangle*, mark size=3pt, only marks, color=impossiblered]
  coordinates {(0.01,4.0278)};
\addlegendentry{adversarial minimum}
\addplot[dashed, gray, domain=0:2.05] {1};
\node[font=\scriptsize, anchor=west, gray] at (axis cs:0.05,1.75)
  {ratio \(=1\): viscosity exactly matches the blow-up rate};
\end{axis}
\end{tikzpicture}
\caption[Suppression ratio against \(\lambda\)]{Suppression ratio against the
self-similar exponent \(\lambda\). Every measured value is above unity ---
viscous damping outpaces the self-similar blow-up rate at every \(\lambda\)
tested, including the adversarial minimum at \(\lambda=0.01\) --- and the
ratio decreases monotonically toward small \(\lambda\) without crossing.
See Observation~\ref{obs:lambda-caveats} before drawing any conclusion from
this figure.}
\label{fig:lambda-sweep}
\end{figure}

\begin{obsv}[How much weight the \(\lambda\) sweep can bear: not much]
\label{obs:lambda-caveats}
Three facts about these ten runs, all read directly from the database, sharply
limit what they establish.

\emph{(i) The runs are five time steps long.} Every row has
\texttt{n\_steps}\(=5\) at \texttt{t\_end}\(=0.5\). These are not evolutions;
they are barely-perturbed initial conditions.

\emph{(ii) The temperature field does not move.} For all nine
CCF and Boussinesq runs, \texttt{T\_max\_final} lies between
\(300.0000000055\) and \(300.0000000064\,\mathrm{K}\): the temperature changed
by about \(6\times10^{-9}\,\mathrm{K}\). Consequently \(A_{\min}\) equals
\(A(300\,\mathrm{K})=3.0801\times10^{-5}\) in all nine, to five significant
figures, and the assertion \(A_{\min}>0\) is not testing the dynamics --- it
is testing the property engine. Only the adversarial run has real thermal
dynamics: its temperature spans \(304.99\) to \(307.49\,\mathrm{K}\), and it
is the only run whose \(A_{\min}\) differs, at \(3.1740\times10^{-5}\).

\emph{(iii) The suppression ratio is largely analytic.} The three Boussinesq
runs report \emph{identical} \(\omega_{\max}=0.9874582480137633\) to sixteen
digits, and identical final temperatures, while their suppression ratios
differ by \(60\%\). The reason is that \(\lambda\) does not enter the
Boussinesq initial condition used --- the dynamics are the same run three
times --- and the suppression ratio is a post-processed comparison of
\(\mu(t)\) against the \emph{analytic} blow-up rate \((T^*-t)^{-(1+\lambda)}\)
evaluated at that \(\lambda\). The monotone descent in
Figure~\ref{fig:lambda-sweep} is therefore in large part a property of the
formula being compared against, not a measurement of the flow.

The honest statement of what M3 shows is: over the specified \(\lambda\) grid,
in short 2D runs, nothing anomalous occurred, and no threshold appeared. It is
consistent with the conjecture. It does not test it.
\end{obsv}

\section{The Route~2 \texorpdfstring{\(\varepsilon\)}{eps} sweep (M2)}
\label{sec:route2-2d}

Four runs at \(\varepsilon\in\{1,0.1,0.01,0.001\}\), \(N=64\),
\(\nu=10^{-3}\), on the CCF \(\lambda=0.6057\) profile.

\begin{center}
\small
\begin{tabular}{@{}l r r r l@{}}
\toprule
\texttt{exp\_id} & \(\varepsilon\) & \(\theta_{\max}\) &
\(\mu_{\mathrm{eff},\min}\) & verdict \\
\midrule
\texttt{EXP-L2-R2-001} & \(1.000\) & \(2.6045\!\times\!10^{-4}\) & \(1.0\!\times\!10^{-3}\) & \textsf{PASS} \\
\texttt{EXP-L2-R2-002} & \(0.100\) & \(2.6045\!\times\!10^{-4}\) & \(1.0\!\times\!10^{-3}\) & \textsf{PASS} \\
\texttt{EXP-L2-R2-003} & \(0.010\) & \(2.6045\!\times\!10^{-4}\) & \(1.0\!\times\!10^{-3}\) & \textsf{PASS} \\
\texttt{EXP-L2-R2-004} & \(0.001\) & \(2.6045\!\times\!10^{-4}\) & \(1.0\!\times\!10^{-3}\) & \textsf{PASS} \\
\bottomrule
\end{tabular}
\end{center}

\medskip
The four \(\theta_{\max}\) values are identical because the source
\(S_\theta=\nu\abs{\nabla u}^2\) does not contain \(\varepsilon\), and the
vorticity evolution is identical because \(\varepsilon\) never enters the
momentum equation at all. This is the numerical face of the structural point
made in Chapter~\ref{ch:routes}: Route~2's auxiliary scalar is diagnostic, and
its regularisation parameter is inert.\footnote{A schema note for anyone
querying the database directly: the Route-2 rows reuse the
\texttt{T\_min\_final} and \texttt{T\_max\_final} columns of the shared
\texttt{experiments} table to store \(\theta\) diagnostics, since Route~2 has
no temperature field. The \texttt{key\_metric} string is authoritative:
\texttt{theta\_max\_final=2.6045e-04}. The column names are a legacy of the
table having been designed for Route~1.}

The recorded \texttt{lps\_margin\_min}\(=0\) is also not a failure: at \(t=0\)
we have \(\theta\equiv0\) and hence \(\varepsilon f(\theta)=0\), so the
minimum over the run is attained at the first step by construction.

\section{The \texorpdfstring{\(\mu(T)\to\nu\)}{mu -> nu} limit in two dimensions (M4)}
\label{sec:mu-limit-2d}

This is the Phase-4 deliverable in 2D: does the regularity margin collapse as
the temperature field becomes uniform? Ten runs, \(N=64\),
\(\bar T=300\,\mathrm{K}\), \(T_0=\bar T+\delta g(x)\) with
\(\norm{g}_{L^\infty}=1\), \(\delta\) descending over three decades. Reference
values: \(A_{\mathrm{ref}}=3.0801337698\times10^{-5}\),
\(\nu_{\mathrm{ref}}=1.5686\times10^{-5}\).

\begin{table}[ht]
\centering
\small
\begin{tabular}{@{}l r r r r@{}}
\toprule
\texttt{exp\_id} & \(\delta\) & \(A_{\min}\) &
\(\varepsilon(\delta)=A_{\min}-A_{\mathrm{ref}}\) &
\(\varepsilon(\delta)/\delta\) \\
\midrule
\texttt{EXP-L2-R1-MULIMIT-001} & \(1.000\) & \(3.0614\!\times\!10^{-5}\) & \(-1.8700\!\times\!10^{-7}\) & \(-1.870\!\times\!10^{-7}\) \\
\texttt{EXP-L2-R1-MULIMIT-002} & \(0.500\) & \(3.0708\!\times\!10^{-5}\) & \(-9.3562\!\times\!10^{-8}\) & \(-1.871\!\times\!10^{-7}\) \\
\texttt{EXP-L2-R1-MULIMIT-003} & \(0.200\) & \(3.0764\!\times\!10^{-5}\) & \(-3.7440\!\times\!10^{-8}\) & \(-1.872\!\times\!10^{-7}\) \\
\texttt{EXP-L2-R1-MULIMIT-004} & \(0.100\) & \(3.0783\!\times\!10^{-5}\) & \(-1.8723\!\times\!10^{-8}\) & \(-1.872\!\times\!10^{-7}\) \\
\texttt{EXP-L2-R1-MULIMIT-005} & \(0.050\) & \(3.0792\!\times\!10^{-5}\) & \(-9.3620\!\times\!10^{-9}\) & \(-1.872\!\times\!10^{-7}\) \\
\texttt{EXP-L2-R1-MULIMIT-006} & \(0.020\) & \(3.0798\!\times\!10^{-5}\) & \(-3.7449\!\times\!10^{-9}\) & \(-1.872\!\times\!10^{-7}\) \\
\texttt{EXP-L2-R1-MULIMIT-007} & \(0.010\) & \(3.0799\!\times\!10^{-5}\) & \(-1.8725\!\times\!10^{-9}\) & \(-1.873\!\times\!10^{-7}\) \\
\texttt{EXP-L2-R1-MULIMIT-008} & \(0.005\) & \(3.0800\!\times\!10^{-5}\) & \(-9.3626\!\times\!10^{-10}\) & \(-1.873\!\times\!10^{-7}\) \\
\texttt{EXP-L2-R1-MULIMIT-009} & \(0.002\) & \(3.0801\!\times\!10^{-5}\) & \(-3.7450\!\times\!10^{-10}\) & \(-1.873\!\times\!10^{-7}\) \\
\texttt{EXP-L2-R1-MULIMIT-010} & \(0.001\) & \(3.0801\!\times\!10^{-5}\) & \(-1.8725\!\times\!10^{-10}\) & \(-1.873\!\times\!10^{-7}\) \\
\bottomrule
\end{tabular}
\caption[The 2D \(\mu(T)\to\nu\) limit]{The M4 sweep. The last column is the
point: \(\varepsilon(\delta)/\delta\) is constant to four significant figures
across three decades. The approach to the Prize limit is exactly first order.
All ten runs \textsf{PASS}.}
\label{tab:mu-limit-2d}
\end{table}

The physical reading is straightforward. Cold spots in the initial temperature
field have slightly lower \(A\) than \(A(\bar T)\), so \(A_{\min}<
A_{\mathrm{ref}}\); as \(\delta\to0\) the cold spots vanish and \(A_{\min}\)
rises to \(A_{\mathrm{ref}}\) from below, at first order in \(\delta\) with
coefficient \(\abs{A'(\bar T)}\). Nothing degenerates. The limit is regular,
not singular.

\begin{obsv}[A column name that misleads]
\label{obs:lps-margin-column}
The quantity tabulated above is stored in a database column named
\texttt{lps\_margin}, and it is \emph{negative} in all ten rows while the
verdict is \textsf{PASS}. A reader who takes the column name at face value
would conclude the programme recorded ten negative regularity margins as
passes. It did not: the quantity is
\(\varepsilon(\delta)=A_{\min}-A_{\mathrm{ref}}\), and its sign is negative
for the reason just given. The \(\lambda\)-sweep and \(\mu\)-limit drivers
share a table; the column name belongs to the former.
\end{obsv}

\section{What Layer~2 established}
\label{sec:layer2-summary}

\begin{center}
\small
\begin{tabular}{@{}l p{0.36\textwidth} p{0.31\textwidth}@{}}
\toprule
Milestone & Claim targeted & Outcome \\
\midrule
M0 & Route-1 2D solver correct; \(\nabla\cdot u=0\) exact & \(64/64\) tests;
  solver is the base of M3 and M4 \\
M2 & Route-2 \(\varepsilon\)-independence & \(4/4\) \textsf{PASS};
  \(\theta_{\max}\) identical across \(\varepsilon\) \\
M3 & No critical \(\lambda_c\) & \(10/10\) \textsf{PASS}; consistent with the
  conjecture, but see Obs.~\ref{obs:lambda-caveats} \\
M4 & \(\mu(T)\to\nu\) limit regular & \(10/10\) \textsf{PASS};
  \(\varepsilon(\delta)\propto\delta\) exactly \\
\bottomrule
\end{tabular}
\end{center}

\medskip
Two of these four are strong. The \(\varepsilon\)-independence result is
strong because it is exact rather than approximate: the same numbers appear at
every \(\varepsilon\) because the same equations were solved. The
\(\mu\)-limit result is strong because the linearity holds to four figures
across three decades, which is not the sort of thing that happens by accident.
The other two are weaker than their \textsf{PASS} verdicts suggest, and
Observation~\ref{obs:lambda-caveats} says why.

None of the four is a theorem, and the programme never claimed otherwise. What
Layer~2 delivered to the analytical work was a validated solver, a set of
adversarial initial conditions, and the habit of logging every number.
